\begin{tikzpicture}[font=\small,>=stealth]
\draw[gray!45,rounded corners=2pt] (0,0) rectangle (14.2,10);
\node at (7.1,9.35) {梯度垂直弦方向：看的是经过中点的等值线};
\node at (7.1,8.5) {$F(x,y)=x^2/4+y^2/3-1,\quad M(1,1),\quad \nabla F(M)=(1/2,2/3)$};
\pgfmathsetmacro{\u}{sqrt(20/21)}
\begin{scope}[shift={(5.2,4.4)},scale=1.4]
\draw[blue!75!black,thick] (0,0) ellipse (2 and 1.7320508076);
\draw[teal!80!black,dashed,thick] (0,0) ellipse (1.5275252317 and 1.3228756555);
\coordinate (A) at ({1-\u},{1+.75*\u});
\coordinate (B) at ({1+\u},{1-.75*\u});
\draw[orange!85!black,very thick] (A)--(B);
\draw[purple!80!black,very thick,->] (1,1)--(1.9,2.2)
 node[right=6pt] {$\nabla F(M)$};
\draw[gray] (1.144,.892)--(1.252,1.036)--(1.108,1.144);
\fill (A) circle(1.5pt) node[above left=5pt] {$A$};
\fill (B) circle(1.5pt) node[below right=5pt] {$B$};
\fill[purple!80!black] (1,1) circle(1.7pt) node[below left=6pt] {$M$};
\end{scope}
\node[text=blue!75!black] at (10.9,5.7) {原曲线：$F=0$};
\node[text=teal!80!black] at (10.9,4.5) {经过 $M$ 的等值线};
\node[text=teal!80!black] at (10.9,3.7) {$F=F(M)=-5/12$};
\node at (7.1,1.3) {$\nabla F(M)\cdot\overrightarrow{AB}=0$，故弦方向与梯度垂直。};
\node at (7.1,.55) {$M$ 不在蓝色原曲线上，所以不能称作“原曲线在 $M$ 处的切线”。};
\end{tikzpicture}
